\documentclass[11pt,fleqn, openany]{book} % Default font size and left-justified equations

%%%%%%%%%%%%%%%%%%%%%%%%%%%%%%%%%%%%%%%%%
% The Legrand Orange Book
% Structural Definitions File
% Version 2.1 (26/09/2018)
%
% Original author:
% Mathias Legrand (legrand.mathias@gmail.com) with modifications by:
% Vel (vel@latextemplates.com)
% 
% This file was downloaded from:
% http://www.LaTeXTemplates.com
%
% License:
% CC BY-NC-SA 3.0 (http://creativecommons.org/licenses/by-nc-sa/3.0/)
%
%%%%%%%%%%%%%%%%%%%%%%%%%%%%%%%%%%%%%%%%%

%----------------------------------------------------------------------------------------
%	VARIOUS REQUIRED PACKAGES AND CONFIGURATIONS
%----------------------------------------------------------------------------------------

\usepackage[table]{xcolor}

\usepackage{graphicx}
\usepackage{tabularx} % Required for including pictures
\usepackage{pgf,tikz,tkz-tab,eurosym,yhmath, stmaryrd}
\usepackage{pgfplots}
\usepackage{mathrsfs}
\usetikzlibrary{patterns}
\usetikzlibrary{trees}
\graphicspath{{../../Pictures/}}
\usepackage{multicol} 


\usepackage[english]{babel} % English language/hyphenation
\usepackage{icomma}
\usepackage{enumitem} % Customize lists
\setlist{nolistsep, nosep, nolistsep} % Reduce spacing between bullet points and numbered lists

\usepackage{booktabs} % Required for nicer horizontal rules in tables

 % Required for specifying colors by name


\definecolor{ocre}{RGB}{243,102,25} % Define the orange color used for highlighting throughout the book

\usepackage{listings}

\definecolor{codegreen}{rgb}{0,0.6,0}
\definecolor{codegray}{rgb}{0.5,0.5,0.5}
\definecolor{codepurple}{rgb}{0.58,0,0.82}
\definecolor{backcolour}{rgb}{0.95,0.95,0.92}

\lstdefinestyle{mystyle}{
    backgroundcolor=\color{backcolour},   
    commentstyle=\color{codegreen},
    keywordstyle=\color{magenta},
    numberstyle=\tiny\color{codegray},
    stringstyle=\color{codepurple},
    basicstyle=\ttfamily\footnotesize,
    breakatwhitespace=false,         
    breaklines=true,                 
    captionpos=b,                    
    keepspaces=true,                 
    numbers=left,                    
    numbersep=5pt,                  
    showspaces=false,                
    showstringspaces=false,
    showtabs=false,                  
    tabsize=2
}

\lstset{style=mystyle}

%----------------------------------------------------------------------------------------
% Paramétrage XSIM
%----------------------------------------------------------------------------------------

\usepackage[no-files]{xsim}


\DeclareExerciseEnvironmentTemplate{myex}{%
    \textbf{%
      \hypertarget{ex:\ExerciseID}{\sffamily{\ensuremath{\blacktriangleright}} Exercice \GetExerciseProperty{counter} \GetExerciseProperty{subtitle} --}
      \hyperlink{sol:\ExerciseID}{Voir le corrigé}%
    }\par
}{\par\smallskip}

\DeclareExerciseEnvironmentTemplate{mysol}{%
    \textbf{%
      \hypertarget{sol:\ExerciseID}{\sffamily{\ensuremath{\blacktriangleright}} Correction \GetExerciseProperty{counter} --}
      \hyperlink{ex:\ExerciseID}{Voir l'énoncé}%
    }\par
}{\par\medskip}

\xsimsetup{
  exercise/template = myex ,
  solution/template = mysol 
}

%Collection exercices

\DeclareExerciseTagging{topic}

\xsimsetup{collect}

%----------------------------------------------------------------------------------------
% SYMBOLES
%----------------------------------------------------------------------------------------

\newcommand\imCMsym[4][\mathord]{%
  \DeclareFontFamily{U} {#2}{}
  \DeclareFontShape{U}{#2}{m}{n}{
    <-6> #25
    <6-7> #26
    <7-8> #27
    <8-9> #28
    <9-10> #29
    <10-12> #210
    <12-> #212}{}
  \DeclareSymbolFont{CM#2} {U} {#2}{m}{n}
  \DeclareMathSymbol{#4}{#1}{CM#2}{#3}
}
\newcommand\alsoimCMsym[4][\mathord]{\DeclareMathSymbol{#4}{#1}{CM#2}{#3}}

\imCMsym{cmmi}{124}{\CMjmath}

\newcommand{\Oij}{(O\,;\,\vec{\imath}\,,\, \vec{\CMjmath} )}
\newcommand{\Oijk}{(O\,;\,\vec{\imath}\,,\, \vec{\CMjmath}\,,\,\vec{k})}

\newcommand\e{\mathrm{e}}
\newcommand\R{\mathbb{R}}
\newcommand\N{\mathbb{N}}


%----------------------------------------------------------------------------------------
%	MARGINS
%----------------------------------------------------------------------------------------

\usepackage{geometry} % Required for adjusting page dimensions and margins

\geometry{
	paper=a4paper, % Paper size, change to letterpaper for US letter size
	top=3cm, % Top margin
	bottom=3cm, % Bottom margin
	left=2cm, % Left margin
	right=2cm, % Right margin
	headheight=14pt, % Header height
	footskip=1.4cm, % Space from the bottom margin to the baseline of the footer
	headsep=10pt, % Space from the top margin to the baseline of the header
	%showframe, % Uncomment to show how the type block is set on the page
}

\setlength{\parindent}{0pt}
\parskip=5pt



%----------------------------------------------------------------------------------------
%	FONTS
%----------------------------------------------------------------------------------------

\usepackage{avant} % Use the Avantgarde font for headings
\usepackage{times} % Use the Times font for headings
\usepackage{mathptmx} % Use the Adobe Times Roman as the default text font together with math symbols from the Sym­bol, Chancery and Com­puter Modern fonts

%\usepackage{microtype} % Slightly tweak font spacing for aesthetics
%\usepackage[utf8]{inputenc} % Required for including letters with accents
\usepackage[T1]{fontenc} % Use 8-bit encoding that has 256 glyphs

%----------------------------------------------------------------------------------------
%	BIBLIOGRAPHY AND INDEX
%----------------------------------------------------------------------------------------

\usepackage[style=numeric,citestyle=numeric,sorting=nyt,sortcites=true,autopunct=true,babel=hyphen,hyperref=true,abbreviate=false,backref=true,backend=biber]{biblatex}
\addbibresource{bibliography.bib} % BibTeX bibliography file
\defbibheading{bibempty}{}

\usepackage{calc} % For simpler calculation - used for spacing the index letter headings correctly
\usepackage{makeidx} % Required to make an index
\makeindex % Tells LaTeX to create the files required for indexing

%----------------------------------------------------------------------------------------
%	MAIN TABLE OF CONTENTS
%----------------------------------------------------------------------------------------

\usepackage{titletoc} % Required for manipulating the table of contents

\contentsmargin{0cm} % Removes the default margin

% Part text styling (this is mostly taken care of in the PART HEADINGS section of this file)
\titlecontents{part}
	[0cm] % Left indentation
	{\addvspace{20pt}\bfseries} % Spacing and font options for parts
	{}
	{}
	{}

% Chapter text styling
\titlecontents{chapter}
	[1.25cm] % Left indentation
	{\addvspace{12pt}\large\sffamily\bfseries} % Spacing and font options for chapters
	{\color{ocre!60}\contentslabel[\Large\thecontentslabel]{1.25cm}\color{ocre}} % Formatting of numbered sections of this type
	{\color{ocre}} % Formatting of numberless sections of this type
	{\color{ocre!60}\normalsize\;\titlerule*[.5pc]{.}\;\thecontentspage} % Formatting of the filler to the right of the heading and the page number

% Section text styling
\titlecontents{section}
	[1.25cm] % Left indentation
	{\addvspace{3pt}\sffamily\bfseries} % Spacing and font options for sections
	{\contentslabel[\thecontentslabel]{1.25cm}} % Formatting of numbered sections of this type
	{} % Formatting of numberless sections of this type
	{\hfill\color{black}\thecontentspage} % Formatting of the filler to the right of the heading and the page number

% Subsection text styling
\titlecontents{subsection}
	[1.25cm] % Left indentation
	{\addvspace{1pt}\sffamily\small} % Spacing and font options for subsections
	{\contentslabel[\thecontentslabel]{1.25cm}} % Formatting of numbered sections of this type
	{} % Formatting of numberless sections of this type
	{\ \titlerule*[.5pc]{.}\;\thecontentspage} % Formatting of the filler to the right of the heading and the page number

% Figure text styling
\titlecontents{figure}
	[1.25cm] % Left indentation
	{\addvspace{1pt}\sffamily\small} % Spacing and font options for figures
	{\thecontentslabel\hspace*{1em}} % Formatting of numbered sections of this type
	{} % Formatting of numberless sections of this type
	{\ \titlerule*[.5pc]{.}\;\thecontentspage} % Formatting of the filler to the right of the heading and the page number

% Table text styling
\titlecontents{table}
	[1.25cm] % Left indentation
	{\addvspace{1pt}\sffamily\small} % Spacing and font options for tables
	{\thecontentslabel\hspace*{1em}} % Formatting of numbered sections of this type
	{} % Formatting of numberless sections of this type
	{\ \titlerule*[.5pc]{.}\;\thecontentspage} % Formatting of the filler to the right of the heading and the page number

%----------------------------------------------------------------------------------------
%	MINI TABLE OF CONTENTS IN PART HEADS
%----------------------------------------------------------------------------------------

% Chapter text styling
\titlecontents{lchapter}
	[0em] % Left indentation
	{\addvspace{15pt}\large\sffamily\bfseries} % Spacing and font options for chapters
	{\color{ocre}\contentslabel[\Large\thecontentslabel]{1.25cm}\color{ocre}} % Chapter number
	{}  
	{\color{ocre}\normalsize\sffamily\bfseries\;\titlerule*[.5pc]{.}\;\thecontentspage} % Page number

% Section text styling
\titlecontents{lsection}
	[0em] % Left indentation
	{\sffamily\small} % Spacing and font options for sections
	{\contentslabel[\thecontentslabel]{1.25cm}} % Section number
	{}
	{}

% Subsection text styling (note these aren't shown by default, display them by searchings this file for tocdepth and reading the commented text)
\titlecontents{lsubsection}
	[.5em] % Left indentation
	{\sffamily\footnotesize} % Spacing and font options for subsections
	{\contentslabel[\thecontentslabel]{1.25cm}}
	{}
	{}

%----------------------------------------------------------------------------------------
%	HEADERS AND FOOTERS
%----------------------------------------------------------------------------------------


\usepackage{fancyhdr} % Required for header and footer configuration

\pagestyle{fancy}
\renewcommand{\chaptermark}[1]{\markboth{\sffamily\normalsize\bfseries\ \thechapter.\ #1}{}} % Chapter text font settings
\renewcommand{\sectionmark}[1]{\markright{\sffamily\normalsize\thesection\hspace{5pt}#1}{}} % Section text font settings
\fancyhf{} \fancyhead[LE,RO]{\sffamily\normalsize\thepage} % Font setting for the page number in the header
\fancyhead[LO]{\rightmark} % Print the nearest section name on the left side of odd pages
\fancyhead[RE]{\leftmark} % Print the current chapter name on the right side of even pages

\fancyfoot[L]{Jason LAPEYRONNIE}
\fancyfoot[R]{\href{http://mathoutils.fr}{http://mathoutils.fr}} % Uncomment to include a footer

\renewcommand{\headrulewidth}{0.5pt} % Thickness of the rule under the header
\renewcommand{\footrulewidth}{0.5pt} % Thickness of the rule under the header

\fancypagestyle{plain}{% Style for when a plain pagestyle is specified
	\fancyhead{}\renewcommand{\headrulewidth}{0pt}%
}

% Removes the header from odd empty pages at the end of chapters
\makeatletter
\renewcommand{\cleardoublepage}{
\clearpage\ifodd\c@page\else
\hbox{}
\vspace*{\fill}
\thispagestyle{empty}
\newpage
\fi}

%----------------------------------------------------------------------------------------
%	THEOREM STYLES
%----------------------------------------------------------------------------------------

\usepackage{amsmath,amsfonts,amssymb,amsthm} % For math equations, theorems, symbols, etc

\newcommand{\intoo}[2]{\mathopen{]}#1\,;#2\mathclose{[}}
\newcommand{\ud}{\mathop{\mathrm{{}d}}\mathopen{}}
\newcommand{\intff}[2]{\mathopen{[}#1\,;#2\mathclose{]}}
\renewcommand{\qedsymbol}{$\blacksquare$}
\newtheorem{notation}{Notation}[section]

% Boxed/framed environments
\newtheoremstyle{ocrenumbox}% Theorem style name
{0pt}% Space above
{0pt}% Space below
{\normalfont}% Body font
{}% Indent amount
{\small\bf\sffamily\color{ocre}}% Theorem head font
{\;:\;}% Punctuation after theorem head
{0.25em}% Space after theorem head
{\small\sffamily\color{ocre}\thmname{#1}\nobreakspace\thmnumber{\@ifnotempty{#1}{}\@upn{#2}}% Theorem text (e.g. Theorem 2.1)
\thmnote{\nobreakspace\the\thm@notefont\sffamily\bfseries\color{black}---\nobreakspace#3}} % Optional theorem note

\newtheoremstyle{blacknumex}% Theorem style name
{5pt}% Space above
{10pt}% Space below
{\normalfont}% Body font
{} % Indent amount
{\small\bf\sffamily}% Theorem head font
{\;:\;}% Punctuation after theorem head
{0.25em}% Space after theorem head
{\small\sffamily{\tiny\ensuremath{\blacksquare}}\nobreakspace\thmname{#1}\nobreakspace\thmnumber{\@ifnotempty{#1}{}\@upn{#2}}% Theorem text (e.g. Theorem 2.1)
\thmnote{\nobreakspace\the\thm@notefont\sffamily\bfseries---\nobreakspace#3}}% Optional theorem note

\newtheoremstyle{blacknumexo}% Theorem style name
{15pt}% Space above
{10pt}% Space below
{\normalfont}% Body font
{} % Indent amount
{\small\bf\sffamily}% Theorem head font
{}% Punctuation after theorem head
{0.5em}% Space after theorem head
{\small\sffamily{\ensuremath{\blacktriangleright}}\nobreakspace\thmname{#1}\nobreakspace\thmnumber{\@ifnotempty{#1}{}\@upn{#2}}% Theorem text (e.g. Theorem 2.1)
\thmnote{\nobreakspace\the\thm@notefont\sffamily\bfseries---\nobreakspace#3} \\}% Optional theorem note



\newtheoremstyle{blacknumbox} % Theorem style name
{0pt}% Space above
{5pt}% Space below
{}% Body font
{}% Indent amount
{\large\bf\sffamily}% Theorem head font
{\;:\;}% Punctuation after theorem head
{0.25em}% Space after theorem head
{\small\sffamily\thmname{#1}\nobreakspace\thmnumber{\@ifnotempty{#1}{}\@upn{#2}}% Theorem text (e.g. Theorem 2.1)
\thmnote{\nobreakspace\the\thm@notefont\sffamily\bfseries---\nobreakspace#3}}% Optional theorem note

% Non-boxed/non-framed environments
\newtheoremstyle{ocrenum}% Theorem style name
{5pt}% Space above
{5pt}% Space below
{\normalfont}% Body font
{}% Indent amount
{\small\bf\sffamily\color{ocre}}% Theorem head font
{\;:\;}% Punctuation after theorem head
{0.25em}% Space after theorem head
{\small\sffamily\color{ocre}\thmname{#1}\nobreakspace\thmnumber{\@ifnotempty{#1}{}\@upn{#2}}% Theorem text (e.g. Theorem 2.1)
\thmnote{\nobreakspace\the\thm@notefont\sffamily\bfseries\color{black}---\nobreakspace#3}} % Optional theorem note
\makeatother

% Defines the theorem text style for each type of theorem to one of the three styles above
\newcounter{dummy} 
\newcounter{thm}
\newcounter{correction}
\newcounter{qst}
\theoremstyle{ocrenumbox}
\newtheorem{theoremeT}[dummy]{Théorème}
\newtheorem{exerciseT}{Propriété}
\newtheorem{principeT}{Principe}
\theoremstyle{blacknumex}
\newtheorem{exampleT}{Exemple}
\theoremstyle{blacknumexo}
\newtheorem{exo}[thm]{Exercice}
\newtheorem{corr}[correction]{Correction}
\newtheorem{quest}[qst]{Question}
\theoremstyle{blacknumbox}
\newtheorem{vocabulary}{Vocabulary}[section]
\newtheorem{definitionT}{Définition}
\newtheorem{corollaryT}[dummy]{Corollary}
\theoremstyle{ocrenum}
\newtheorem{proofT}[dummy]{Démonstration}


%----------------------------------------------------------------------------------------
%	DEFINITION OF COLORED BOXES
%----------------------------------------------------------------------------------------

\RequirePackage[framemethod=default]{mdframed} % Required for creating the theorem, definition, exercise and corollary boxes

% Theorem box
\newmdenv[skipabove=7pt,
skipbelow=7pt,
backgroundcolor=black!5,
linecolor=ocre,
innerleftmargin=5pt,
innerrightmargin=5pt,
innertopmargin=10pt,
leftmargin=0cm,
rightmargin=0cm,
innerbottommargin=5pt]{tBox}

%Proposition box	  
\newmdenv[skipabove=7pt,
skipbelow=7pt,
rightline=false,
leftline=true,
topline=false,
bottomline=false,
backgroundcolor=ocre!10,
linecolor=ocre,
innerleftmargin=5pt,
innerrightmargin=5pt,
innertopmargin=10pt,
innerbottommargin=3pt,
leftmargin=0cm,
rightmargin=0cm,
linewidth=4pt]{eBox}	

% Definition box
\newmdenv[skipabove=7pt,
backgroundcolor=ocre!4,
skipbelow=7pt,
rightline=false,
leftline=true,
topline=false,
bottomline=false,
linecolor=ocre,
innerleftmargin=5pt,
innerrightmargin=5pt,
innertopmargin=10pt,
leftmargin=0cm,
rightmargin=0cm,
linewidth=4pt,
innerbottommargin=5pt]{dBox}	

% Corollary box
\newmdenv[skipabove=7pt,
skipbelow=7pt,
rightline=false,
leftline=true,
topline=false,
bottomline=false,
linecolor=gray,
backgroundcolor=black!5,
innerleftmargin=5pt,
innerrightmargin=5pt,
innertopmargin=5pt,
leftmargin=0cm,
rightmargin=0cm,
linewidth=4pt,
innerbottommargin=5pt]{cBox}

\newmdenv[skipabove=7pt,
skipbelow=7pt,
backgroundcolor=black!5,
innerleftmargin=5pt,
topline=false,
bottomline=false,
rightline=false,
leftline=false,
innerrightmargin=5pt,
innertopmargin=5pt,
leftmargin=0cm,
rightmargin=0cm,
innerbottommargin=5pt]{xBox}

% Creates an environment for each type of theorem and assigns it a theorem text style from the "Theorem Styles" section above and a colored box from above
\newenvironment{theorem}{\begin{tBox}\begin{theoremeT}}{\end{theoremeT}\end{tBox}}

\newenvironment{exo2}{\noindent \begin{exo}\item\relax \noindent \begin{eBox}\item\relax}{\end{eBox}\end{exo}}


\newenvironment{proposition}{\begin{eBox}\begin{exerciseT}}{\hfill{\color{ocre}}\end{exerciseT}\end{eBox}}		

\newenvironment{principe}{\begin{eBox}\begin{principeT}}{\hfill{\color{ocre}}\end{principeT}\end{eBox}}	
		  
\newenvironment{definition}{\begin{dBox}\begin{definitionT}}{\end{definitionT}\end{dBox}}	

\newenvironment{example}{\begin{xBox}\begin{exampleT}}{\hfill{\tiny\ensuremath{\blacksquare}}\end{exampleT}\end{xBox}}

\newenvironment{demonstration}{\begin{proofT}}{\hfill{\tiny\ensuremath{\square}}\end{proofT}}		
\newenvironment{corollary}{\begin{cBox}\begin{corollaryT}}{\end{corollaryT}\end{cBox}}	

%----------------------------------------------------------------------------------------
%	REMARK ENVIRONMENT
%----------------------------------------------------------------------------------------

\newenvironment{remark}{\par\vspace{5pt}\small % Vertical white space above the remark and smaller font size
\begin{list}{}{
\leftmargin=25pt % Indentation on the left
\rightmargin=15pt}\item\ignorespaces % Indentation on the right
\makebox[-2.5pt]{\begin{tikzpicture}[overlay]
\node[draw=ocre!60,line width=1pt,circle,fill=ocre!25,font=\sffamily\bfseries,inner sep=2pt,outer sep=0pt] at (-15pt,0pt){\textcolor{ocre}{R}};\end{tikzpicture}} % Orange R in a circle
\advance\baselineskip -1pt}{\end{list}\vskip5pt} % Tighter line spacing and white space after remark

%----------------------------------------------------------------------------------------
%	SECTION NUMBERING IN THE MARGIN
%----------------------------------------------------------------------------------------

\makeatletter
\renewcommand{\@seccntformat}[1]{\llap{\textcolor{ocre}{\csname the#1\endcsname}\hspace{1em}}}                    
\renewcommand{\section}{\@startsection{section}{1}{\z@}
{-4ex \@plus -1ex \@minus -.4ex}
{1ex \@plus.2ex }
{\normalfont\LARGE\sffamily\bfseries}}
\renewcommand{\subsection}{\@startsection {subsection}{2}{\z@}
{-3ex \@plus -0.1ex \@minus -.4ex}
{0.5ex \@plus.2ex }
{\normalfont\sffamily\bfseries}}
\renewcommand{\subsubsection}{\@startsection {subsubsection}{3}{\z@}
{-2ex \@plus -0.1ex \@minus -.2ex}
{.2ex \@plus.2ex }
{\normalfont\small\sffamily\bfseries}}                        
\renewcommand\paragraph{\@startsection{paragraph}{4}{\z@}
{-2ex \@plus-.2ex \@minus .2ex}
{.1ex}
{\normalfont\small\sffamily\bfseries}}

%----------------------------------------------------------------------------------------
%	PART HEADINGS
%----------------------------------------------------------------------------------------

% Numbered part in the table of contents
\newcommand{\@mypartnumtocformat}[2]{%
	\setlength\fboxsep{0pt}%
	\noindent\colorbox{ocre!20}{\strut\parbox[c][.7cm]{\ecart}{\color{ocre!70}\Large\sffamily\bfseries\centering#1}}\hskip\esp\colorbox{ocre!40}{\strut\parbox[c][.7cm]{\linewidth-\ecart-\esp}{\Large\sffamily\centering#2}}%
}

% Unnumbered part in the table of contents
\newcommand{\@myparttocformat}[1]{%
	\setlength\fboxsep{0pt}%
	\noindent\colorbox{ocre!40}{\strut\parbox[c][.7cm]{\linewidth}{\Large\sffamily\centering#1}}%
}

\newlength\esp
\setlength\esp{4pt}
\newlength\ecart
\setlength\ecart{1.2cm-\esp}
\newcommand{\thepartimage}{}%
\newcommand{\partimage}[1]{\renewcommand{\thepartimage}{#1}}%
\def\@part[#1]#2{%
\ifnum \c@secnumdepth >-2\relax%
\refstepcounter{part}%
\addcontentsline{toc}{part}{\texorpdfstring{\protect\@mypartnumtocformat{\thepart}{#1}}{\partname~\thepart\ ---\ #1}}
\else%
\addcontentsline{toc}{part}{\texorpdfstring{\protect\@myparttocformat{#1}}{#1}}%
\fi%
\startcontents%
\markboth{}{}%
{\thispagestyle{empty}%
\begin{tikzpicture}[remember picture,overlay]%
\node at (current page.north west){\begin{tikzpicture}[remember picture,overlay]%	
\fill[ocre!20](0cm,0cm) rectangle (\paperwidth,-\paperheight);
\node[anchor=north] at (4cm,-3.25cm){\color{ocre!40}\fontsize{220}{100}\sffamily\bfseries\thepart}; 
\node[anchor=south east] at (\paperwidth-1cm,-\paperheight+1cm){\parbox[t][][t]{8.5cm}{
\printcontents{l}{0}{\setcounter{tocdepth}{1}}% The depth to which the Part mini table of contents displays headings; 0 for chapters only, 1 for chapters and sections and 2 for chapters, sections and subsections
}};
\node[anchor=north east] at (\paperwidth-1.5cm,-3.25cm){\parbox[t][][t]{15cm}{\strut\raggedleft\color{white}\fontsize{30}{30}\sffamily\bfseries#2}};
\end{tikzpicture}};
\end{tikzpicture}}%
\@endpart}
\def\@spart#1{%
\startcontents%
\phantomsection
{\thispagestyle{empty}%
\begin{tikzpicture}[remember picture,overlay]%
\node at (current page.north west){\begin{tikzpicture}[remember picture,overlay]%	
\fill[ocre!20](0cm,0cm) rectangle (\paperwidth,-\paperheight);
\node[anchor=north east] at (\paperwidth-1.5cm,-3.25cm){\parbox[t][][t]{15cm}{\strut\raggedleft\color{white}\fontsize{30}{30}\sffamily\bfseries#1}};
\end{tikzpicture}};
\end{tikzpicture}}
\addcontentsline{toc}{part}{\texorpdfstring{%
\setlength\fboxsep{0pt}%
\noindent\protect\colorbox{ocre!40}{\strut\protect\parbox[c][.7cm]{\linewidth}{\Large\sffamily\protect\centering #1\quad\mbox{}}}}{#1}}%
\@endpart}
\def\@endpart{\vfil\newpage
\if@twoside
\if@openright
\null
\thispagestyle{empty}%
\newpage
\fi
\fi
\if@tempswa
\twocolumn
\fi}

%----------------------------------------------------------------------------------------
%	CHAPTER HEADINGS
%----------------------------------------------------------------------------------------

% A switch to conditionally include a picture, implemented by Christian Hupfer
\newif\ifusechapterimage
\usechapterimagetrue
\newcommand{\thechapterimage}{}%
\newcommand{\chapterimage}[1]{\ifusechapterimage\renewcommand{\thechapterimage}{#1}\fi}%
\newcommand{\autodot}{.}
\def\@makechapterhead#1{%
{\parindent \z@ \raggedright \normalfont
\ifnum \c@secnumdepth >\m@ne
\if@mainmatter
\begin{tikzpicture}[remember picture,overlay]
\node at (current page.north west)
{\begin{tikzpicture}[remember picture,overlay]
\node[anchor=north west,inner sep=0pt] at (0,0) {\ifusechapterimage\includegraphics[width=\paperwidth]{\thechapterimage}\fi};
\draw[anchor=west] (\Gm@lmargin,-3cm) node [line width=2pt,rounded corners=15pt,draw=ocre,fill=white,fill opacity=0.5,inner sep=15pt]{\strut\makebox[22cm]{}};
\draw[anchor=west] (\Gm@lmargin+.3cm,-3cm) node {\huge\sffamily\bfseries\color{black}\thechapter\autodot~#1\strut};
\end{tikzpicture}};
\end{tikzpicture}
\else
\begin{tikzpicture}[remember picture,overlay]
\node at (current page.north west)
{\begin{tikzpicture}[remember picture,overlay]
\node[anchor=north west,inner sep=0pt] at (0,0) {\ifusechapterimage\includegraphics[width=\paperwidth]{\thechapterimage}\fi};
\draw[anchor=west] (\Gm@lmargin,-3cm) node [line width=2pt,rounded corners=15pt,draw=ocre,fill=white,fill opacity=0.5,inner sep=15pt]{\strut\makebox[22cm]{}};
\draw[anchor=west] (\Gm@lmargin+.3cm,-3cm) node {\huge\sffamily\bfseries\color{black}#1\strut};
\end{tikzpicture}};
\end{tikzpicture}
\fi\fi\par\vspace*{50\p@}}}

%-------------------------------------------

\def\@makeschapterhead#1{%
\begin{tikzpicture}[remember picture,overlay]
\node at (current page.north west)
{\begin{tikzpicture}[remember picture,overlay]
\node[anchor=north west,inner sep=0pt] at (0,0) {\ifusechapterimage\includegraphics[width=\paperwidth]{\thechapterimage}\fi};
\draw[anchor=west] (\Gm@lmargin,-3cm) node [line width=2pt,rounded corners=15pt,draw=ocre,fill=white,fill opacity=0.5,inner sep=15pt]{\strut\makebox[22cm]{}};
\draw[anchor=west] (\Gm@lmargin+.3cm,-3cm) node {\huge\sffamily\bfseries\color{black}#1\strut};
\end{tikzpicture}};
\end{tikzpicture}
\par\vspace*{50\p@}}
\makeatother

%----------------------------------------------------------------------------------------
%	LINKS
%----------------------------------------------------------------------------------------

\usepackage{hyperref}
\hypersetup{hidelinks,backref=true,pagebackref=true,hyperindex=true,colorlinks=false,breaklinks=true,urlcolor=ocre,bookmarks=true,bookmarksopen=false}

\usepackage{bookmark}
\bookmarksetup{
open,
numbered,
addtohook={%
\ifnum\bookmarkget{level}=0 % chapter
\bookmarksetup{bold}%
\fi
\ifnum\bookmarkget{level}=-1 % part
\bookmarksetup{color=ocre,bold}%
\fi
}
}

\renewcommand*\thesection{\arabic{section}}

\newcommand*{\coord}[3]{% 
  \ensuremath{\overrightarrow{#1}\, 
    \begin{pmatrix} 
      #2\\ 
      #3 
    \end{pmatrix}}}
    
  \newcommand*{\coordb}[2]{% 
  \ensuremath{ 
    \begin{pmatrix} 
      #1\\ 
      #2 
    \end{pmatrix}}}

\newcommand*{\coorde}[4]{% 
  \renewcommand{\arraystretch}{1}\ensuremath{\overrightarrow{#1}\, 
    \begin{pmatrix} 
      #2\\ 
      #3 \\
      #4
    \end{pmatrix}}}    
  \newcommand*{\coordbe}[3]{% 
 \renewcommand{\arraystretch}{1} \ensuremath{ 
    \begin{pmatrix} 
      #1\\ 
      #2 \\
      #3
    \end{pmatrix}}}  
    
\newcommand{\Card}{\mathrm{Card}}



\begin{document}

\chapterimage{../../Pictures/background}


\chapter{Cours : Orthogonalité dans l'espace}


\section{Produit scalaire de deux vecteurs}


\subsection{Définition du produit scalaire}


\begin{definition}Soient $\vec u$ et $\vec v$ deux vecteurs de l'espace. On considère des points $A$, $B$ et $C$ tels que $\vec u = \overrightarrow{AB}$ et $\vec v= \overrightarrow{AC}$. L'angle non orienté entre les vecteurs $\vec u$ et $\vec v$, noté $(\vec u;\vec v)$ est  l'angle $\widehat{BAC}$, vu dans le plan $(ABC)$.\end{definition}

\begin{definition}Le produit scalaire de $\vec u$ et $\vec v$ est le \textbf{réel} notée $\vec u \cdot \vec v$ et qui vaut
\begin{itemize}
\item 0 si $\vec u$ ou $\vec v$ vaut $\overrightarrow{0}$ ;
\item $\lvert\lvert\vec u\rvert\rvert \times \lvert\lvert\vec v\rvert\rvert \times \cos (\vec u ; \vec v)$ sinon.
\end{itemize}\end{definition}

\textbf{Rappel de certaines valeurs remarquables}

\begin{center}
\renewcommand{\arraystretch}{2.5}
\noindent \begin{tabularx}{0.9\linewidth}{|X|XXXXXX|}
\hline
Degré & 0 & 30 & 45 & 60 & 90 & 180 \\
\hline
Radians &  &  &  &  &  &  \\
\hline
Cosinus &  &  &  &  &  & \\
\hline
Sinus &  &  &  &  & & \\
\hline
\end{tabularx}
\end{center}




\begin{example} Dans un cube $ABCDEFGH$ de côté 1, calculer le produit scalaire $\overrightarrow{AD} \cdot \overrightarrow{BG}$.


\includegraphics[scale=0.7]{cube4}


\end{example}

\begin{definition}On dit que deux vecteurs $\vec u$ et $\vec v$ sont orthogonaux si $\vec u \cdot \vec v=0$.
\end{definition}

Le vecteur $\vec 0$ est en particulier orthogonal à tous les autres vecteurs.

\newpage
\subsection{Propriétés du produit scalaire}

\begin{proposition}Soit $\vec u$, $\vec v$ et $\vec w$ trois vecteurs de l'espace, $k$ et $k'$ deux réels.
\begin{itemize}
\item $\vec u \cdot \vec u = \lvert\lvert \vec u\rvert\rvert^2$. En particulier, $\lvert\lvert\vec u\rvert\rvert=\sqrt{\vec u \cdot \vec u}$.
\item $\vec u \cdot \vec v = \vec v \cdot \vec u$, le produit scalaire est \textbf{symétrique}.
\item $\vec u \cdot ( k \vec v + k' \vec w)= k (\vec u \cdot \vec v) + k' (\vec u \cdot \vec w)$ et $( k \vec v + k' \vec w)\cdot \vec u= k (\vec v \cdot \vec u) + k' (\vec w \cdot \vec u)$. \\Le produit scalaire est \textbf{bilinéaire}.
\end{itemize}\end{proposition}

\begin{example} Soit $\vec u$, $\vec v$ et $\vec w$ trois vecteurs tels que $\vec u \cdot \vec v = 3$, $\vec v \cdot \vec w=5$, $\vec u \cdot \vec w = -1$ et $\lvert\lvert\vec u\rvert\rvert=4$. On a
\[ (\vec u + 2\vec v) \cdot (-3\vec u + 4 \vec w)=\]
On remplace alors les valeurs par celle de l'énoncé en rappelant que $\vec u \cdot \vec u = \lvert\lvert\vec u\rvert\rvert^2$.
\[ (\vec u + 2\vec v) \cdot (-3\vec u + 4 \vec w)=\]\vspace{-0.5cm}\end{example}

\subsection{Formules de polarisation}

\begin{proposition}Soit $\vec u$ et $\vec v$ deux vecteurs de l'espace.
\begin{itemize}
\item $(\vec{u}+\vec{v})^2=\lvert\lvert\vec{u}\rvert\rvert^2+2\vec{u}\cdot\vec{v}+\lvert\lvert\vec{v}\rvert\rvert^2$ ;
\item $(\vec{u}-\vec{v})^2=\lvert\lvert\vec{u}\rvert\rvert^2-2\vec{u}\cdot\vec{v}+\lvert\lvert\vec{v}\rvert\rvert^2$ ;
\item $(\vec{u}+\vec{v})\cdot (\vec{u}-\vec{v})=\lvert\lvert\vec{u}\rvert\rvert^2-\lvert\lvert\vec{v}\rvert\rvert^2$.
\end{itemize}\end{proposition}

\begin{demonstration} On utilise la bilinéarité du produit scalaire. On a\[(\vec{u}+\vec{v})^2 =  (\vec{u}+\vec{v}) \cdot (\vec{u}+\vec{v})= \vec{u}\cdot\vec{u}+\vec{u}\cdot \vec{v}+\vec{v}\cdot\vec{u}+\vec{v}\cdot \vec{v}=\lvert\lvert\vec{u}\rvert\rvert^2+2\vec{u}\cdot\vec{v}+\lvert\lvert\vec{v}\rvert\rvert^2.\]
\[(\vec{u}-\vec{v})^2 =  (\vec{u}-\vec{v}) \cdot (\vec{u}-\vec{v})= \vec{u}\cdot\vec{u}-\vec{u}\cdot \vec{v}-\vec{v}\cdot\vec{u}+\vec{v}\cdot \vec{v}=\lvert\lvert\vec{u}\rvert\rvert^2-2\vec{u}\cdot\vec{v}+\lvert\lvert\vec{v}\rvert\rvert^2.\]
\[(\vec{u}+\vec{v})\cdot (\vec{u}-\vec{v})= \vec{u}\cdot\vec{u}-\vec{u}\cdot \vec{v}+\vec{v}\cdot\vec{u}-\vec{v}\cdot \vec{v}=\lvert\lvert\vec{u}\rvert\rvert^2-\lvert\lvert\vec{v}\rvert\rvert^2.\]
\end{demonstration}

\begin{proposition}Soit $\vec{u}$ et $\vec{v}$ deux vecteurs de l'espace.
\begin{itemize}
\item $\vec{u} \cdot \vec{v}= \dfrac{1}{2}(\lvert\lvert\vec{u}+\vec{v}\rvert\rvert^2-\lvert\lvert\vec{u}\rvert\rvert^2-\lvert\lvert\vec{v}\rvert\rvert^2)$ ;
\item $\vec{u} \cdot \vec{v}= -\dfrac{1}{2}(\lvert\lvert\vec{u}-\vec{v}\rvert\rvert^2-\lvert\lvert\vec{u}\rvert\rvert^2-\lvert\lvert\vec{v}\rvert\rvert^2)$ ;
\item $\vec{u} \cdot \vec{v} = \dfrac{1}{4}( \lvert\lvert\vec{u}+\vec{v}\rvert\rvert^2-\lvert\lvert\vec{u}-\vec{v}\rvert\rvert^2 )$.
\end{itemize}\end{proposition}

\begin{demonstration} Soit $\vec{u}$ et $\vec{v}$ deux vecteurs du plan. D'après la propriété précédente, on a 
\[\lvert\lvert\vec{u}+\vec{v}\rvert\rvert^2=\lvert\lvert\vec{u}\rvert\rvert^2+2\vec{u}\cdot\vec{v}+\lvert\lvert\vec{v}\rvert\rvert^2\quad\text{et}\quad \lvert\lvert\vec{u}-\vec{v}\rvert\rvert^2=\lvert\lvert\vec{u}\rvert\rvert^2-2\vec{u}\cdot\vec{v}+\lvert\lvert\vec{v}\rvert\rvert^2.\]
En isolant $\vec{u}\cdot \vec{v}$ dans ces deux expressions, on retrouve les deux premiers points. De plus, en soustrayant ces deux égalités, on trouve que
\[\lvert\lvert\vec{u}+\vec{v}\rvert\rvert^2-\lvert\lvert\vec{u}-\vec{v}\rvert\rvert^2=4\vec{u}\cdot\vec{v}.\]
Il suffit de diviser par 4 pour retrouver la dernière égalité recherchée.\end{demonstration}

\begin{example}Soit $A$, $B$ et $C$ trois points de l'espace tels que $AB=5$, $BC=7$ et $AC=8$. On a
\[\overrightarrow{AB}\cdot \overrightarrow{AC}=\]
Or, $\overrightarrow{AB}-\overrightarrow{AC}=\overrightarrow{AB}+\overrightarrow{CA}=\overrightarrow{CB}$, d'où
\[\overrightarrow{AB}\cdot \overrightarrow{AC}=\]\vspace{-0.5cm}\end{example}


\section{Base orthonormée}

\begin{definition}Soit $\Oijk$ un repère de l'espace.
\begin{itemize}
\item On dit que la base $(\vec{\imath}, \vec{\CMjmath}, \vec k)$ est orthonormée si
\begin{itemize}
\item $\vec{\imath} \cdot \vec{\CMjmath} = \vec{\imath} \cdot \vec k = \vec{\CMjmath} \cdot \vec k = 0$ ;
\item $\lvert\lvert\vec{\imath}\rvert\rvert=\lvert\lvert\vec{\CMjmath}\rvert\rvert=\lvert\lvert\vec k\rvert\rvert = 1$.
\end{itemize}
\item On dit alors que le repère $\Oijk$ est un repère orthonormé.
\end{itemize}\end{definition}

Une famille orthonormée de trois vecteurs forme forcément une base de l'espace.

\begin{example} Si on considère un cube $ABCDEFGH$ de côté 1, le repère $(A;\overrightarrow{AB},\overrightarrow{AD}, \overrightarrow{AE})$ est un repère orthonormé de l'espace.\end{example}

\begin{proposition}On se place dans un repère orthonormé $\Oijk$. 

On considère deux vecteurs \renewcommand{\arraystretch}{1}\coorde{u}{x}{y}{z} et \renewcommand{\arraystretch}{1}\coorde{v}{x'}{y'}{z'}. Alors, $\vec u \cdot \vec v =xx'+yy'+zz'$.\end{proposition}

\begin{demonstration} Il suffit de revenir à la définition de coordonnées. On a
\[ \vec u \cdot \vec v = (x\vec{\imath} + y \vec{\CMjmath} + z \vec k) \cdot (x' \vec{\imath} + y' \vec{\CMjmath} + z' \vec k).\]
On développe alors
\[ \vec u \cdot \vec v = xx' \vec{\imath} \cdot \vec{\imath} +  xy' \vec{\imath} \cdot \vec{\CMjmath} +  xz' \vec{\imath} \cdot \vec k +  yx' \vec{\CMjmath} \cdot \vec{\imath} +  yy' \vec{\CMjmath} \cdot \vec{\CMjmath} +  yz' \vec{\CMjmath} \cdot \vec k +  zx' \vec k \cdot \vec{\imath} +  zy' \vec k \cdot \vec{\CMjmath} +  zz' \vec k \cdot \vec k .\]
Or, la base $(\vec{\imath}, \vec{\CMjmath}, \vec k)$ est orthonormé, les seuls produits scalaires non nuls sont $\vec{\imath} \cdot \vec{\imath}$, $\vec{\CMjmath} \cdot \vec{\CMjmath}$, et $\vec k \cdot \vec k$ qui valent 1. Ainsi, 
\[ \vec u \cdot \vec v =xx'+yy'+zz'.\]\vspace{-0.5cm}\end{demonstration}

\begin{example} Dans un repère orthonormé $\Oijk$, on considère les vecteurs \renewcommand{\arraystretch}{1}\coorde{u}{3}{-1}{2} et \renewcommand{\arraystretch}{1}\coorde{v}{5}{3}{-6}.
\[ \vec u \cdot \vec v = \]
Les vecteurs $\vec u$ et $\vec v$ sont orthogonaux.\end{example}

\begin{proposition} On se place dans un repère orthonormé $\Oijk$. Soit \renewcommand{\arraystretch}{1}\coorde{u}{x}{y}{z} un vecteur de l'espace. 

Alors $ \lvert\lvert\vec u\rvert\rvert = \sqrt{ \vec u \cdot \vec u} =  \sqrt{x^2+y^2+z^2}$.

En particulier, si $A(x_A,y_A,z_A)$ et $B(x_B, y_B, z_B)$ sont deux points de l'espace, alors
\[ AB = \sqrt{(x_B-x_A)^2+(y_B-y_A)^2+(z_B-z_A)^2} .\]\end{proposition}

\begin{example} On se place dans un repère orthonormé $\Oijk$. Soit $A(1;2;5)$ et $B(3;3;3)$. On a alors
\[ AB = \]\end{example}



\section{Orthogonalité}

\subsection{Droites orthogonales}

\begin{definition}Soit $(d)$ et $(d')$ deux droites de l'espace. On dit que $(d)$ et $(d')$ sont orthogonales si les parallèles à ces deux droites passant par un même point sont perpendiculaires.\end{definition}

\begin{example} On considère un cube $ABCDEFGH$.

\begin{center}
\includegraphics[scale=1]{cube}
\end{center}

Les droites $(AB)$ et $(CG)$ sont orthogonales. En effet, la parallèle à $(CG)$ passant par $B$ est la droite $(BF)$ qui est perpendiculaire à la droite $(AB)$.\end{example}
\newpage
\begin{proposition}Deux droites $(d)$ et $(d')$, dirigées respectivement par les vecteurs $\vec u$ et $\vec v$, sont orthogonales si et seulement si $\vec u \cdot \vec v = 0$.\end{proposition}

\begin{example}On se place dans un repère orthonormé $(O,\vec{\imath}, \vec{\CMjmath}, \vec k)$. On considère les droites $(d_1)$ et $(d_2)$ définies par les représentations paramétriques suivantes.

\[ (d_1)\, ; \, \left\{ \renewcommand{\arraystretch}{1}\begin{array}{l}x=1-2t\\y=3+t\\z=2-t\end{array}\right. \quad \text{et} \quad (d_2)\, ; \, \left\{ \renewcommand{\arraystretch}{1}\begin{array}{l}x=5-3t\\y=3-2t\\z=2+4t\end{array}\right.\]

La droite $(d_1)$ est dirigée par le vecteur $\qquad\qquad$ et la droite $(d_2)$  par le vecteur 
\vskip20pt

Le repère étant orthonormé, on peut calculer le produit scalaire de ces deux vecteurs comme suit :
\[ \vec u \cdot \vec v \]
Les vecteurs $\vec u$ et $\vec v$ sont orthogonaux. Les droites $(d_1)$ et $(d_2)$ sont donc orthogonales.\end{example}

\subsection{Droite orthogonale à un plan}

\begin{definition}Soit $(d)$ une droite et $\mathcal{P}$ un plan de l'espace. On dit que $(d)$ est orthogonale au plan $\mathcal{P}$ si elle est orthogonale à toute droite contenue dans le plan $\mathcal{P}$.\end{definition}

\begin{proposition}[Théorème de la porte]Soit $(d)$ une droite de vecteur directeur $\vec u$ et $\mathcal{P}$ un plan de l'espace dirigé par les vecteurs $\vec v_1$ et $\vec v_2$. La droite $(d)$ est orthogonale au plan $\mathcal{P}$ si et seulement si $\vec u \cdot \vec v_1 = \vec u \cdot \vec v_2 = 0$.\end{proposition}

Il suffit donc de montrer qu'une droite est orthogonale à deux droites sécantes d'un plan pour montrer qu'elle est en fait orthogonale à toute droites de ce plan.


\begin{example} On se place dans un repère orthonormé $\Oijk$. 

On considère les vecteurs \renewcommand{\arraystretch}{1}\coorde{v_1}{1}{2}{3} et \renewcommand{\arraystretch}{1}\coorde{v_2}{1}{0}{1} ainsi que deux points $A(2;5;2)$ et $B(5;8;-1)$.

On note $\mathcal{P}$ le plan passant par le point $O$ et dirigé par les vecteurs $\vec v_1$ et $\vec v_2$ (ces vecteurs n'étant pas colinéaires, on définit bien ainsi un plan). La droite $(AB)$ est orthogonale au plan $\mathcal{P}$. En effet,

\vskip100pt

\end{example}

\newpage

\subsection{Vecteur normal à un plan}

\begin{definition}Soit $\mathcal{P}$ un plan et $\vec n$ un vecteur non nul.
On dit que $\vec n$ est un vecteur normal au plan $\mathcal{P}$ s'il est orthogonal à tout vecteur directeur du plan $\mathcal{P}$.\end{definition}

De la même manière que ce que l'on faisait avec les droites orthogonales à un plan, pour montrer qu'un vecteur est normal à un plan, il suffit de montrer qu'il est orthogonal à deux vecteurs non colinéaires de ce plan.

\begin{example}On se place dans un repère orthonormé $\Oijk$. On considère le plan $\mathcal{P}$ passant par $O$ et dirigé par les vecteurs \renewcommand{\arraystretch}{1}\coorde{v_1}{3}{2}{6} et \renewcommand{\arraystretch}{1}\coorde{v_2}{1}{-3}{2}. Le vecteur \renewcommand{\arraystretch}{1}\coorde{u}{2}{0}{-1} est  normal au plan $\mathcal{P}$. En effet,
\begin{itemize}
\item Les vecteurs $\vec{v_1}$ et $\vec{v}_2$ ne sont pas colinéaires. Ils définissent donc bien un plan.
\item Puisque l'on est dans un repère orthonormé, il est possible de calculer le produit scalaire à l'aide des coordonnées,
\item $ \vec u \cdot \vec v_1 = $
\item $\vec u \cdot \vec v_2 = $
\end{itemize}
Ainsi, $\vec u$ est orthogonal à $\vec v_1$ et $\vec v_2$. $\vec u$ est donc un vecteur normal au plan $\mathcal{P}$.\end{example}

\begin{proposition}Deux plans sont parallèles si et seulement si un vecteur normal au premier est normal au second.\end{proposition}

\begin{example}On se place dans un repère orthonormé $\Oijk$. 

On considère le plan $\mathcal{P}_2$ passant par $A(2;5;9)$, dirigé par \renewcommand{\arraystretch}{1}\coorde{w_1}{5}{4}{10} et \renewcommand{\arraystretch}{1}\coorde{w_2}{-2}{7}{-4}. 

Le vecteur \renewcommand{\arraystretch}{1}\coorde{u}{2}{0}{-1}, qui était normal au plan $\mathcal{P}$ de l'exemple précédent, est aussi normal au plan $\mathcal{P}_2$.\\ En effet, les vecteurs $\vec{w_1}$ et $\vec{w}_2$ ne sont pas colinéaires et définissent donc bien un plan.
De plus, puisque l'on est dans un repère orthonormé, il est possible de calculer le produit scalaire à l'aide des coordonnées,
\begin{itemize}\item $ \vec u \cdot \vec w_1 = 2 \times 5 + 0 \times 4 + (-1) \times 10 = 0$ ;
\item $\vec u \cdot \vec w_2 = 2 \times (-2) + 0 \times 7 + (-1) \times (-4) = 0$.
\end{itemize}
Ainsi, $\vec u$ est orthogonal à $\vec w_1$ et $\vec w_2$. $\vec u$ est donc un vecteur normal au plan $\mathcal{P}_2$. Les plans $\mathcal{P}$ et $\mathcal{P}_2$ admettent un même vecteur normal, ils sont donc parallèles.\end{example}


\begin{definition}Deux plans sont perpendiculaires si le premier plan contient une droite orthogonale au second plan.\end{definition}

Le concept de plans perpendiculaires peut être trompeur. Par exemple, deux plans perpendiculaires peuvent contenir des droites parallèles. En revanche, il est possible de caractériser la perpendicularité de deux plans en utilisant leurs vecteurs normaux.

\begin{proposition} Deux plans sont perpendiculaires si un vecteur normal à l'un est orthogonal à un vecteur normal du second.
\end{proposition}


\newpage

\section{Equation cartésienne d'un plan}

Dans toute cette partie et dans la suivante, l'espace est muni d'un repère orthonormé $\Oijk$.

\subsection{Equation cartésienne}

\begin{proposition}Soit $\vec n$ un vecteur de l'espace et $A$ un point de l'espace. L'ensemble des points $M$ tel que $\overrightarrow{AM}\cdot \vec{n}=0$ est le plan passant par $A$ admettant le vecteur $\vec n$ comme vecteur normal.

Réciproquement, soit $\mathcal{P}$ un plan de l'espace, $A$ un point de $\mathcal{P}$ et $\vec n$ un vecteur normal à $\mathcal{P}$. $\mathcal{P}$ est l'ensemble des points $M$ tels que $\overrightarrow{AM} \cdot \vec n=0$.\end{proposition}

Il est donc possible de décrire un plan à l'aide d'un point et d'un vecteur normal.

\begin{proposition}Soit $A (x_A,y_A,z_A)$ un point de l'espace et $\renewcommand{\arraystretch}{1}\coorde{n}{a}{b}{c}$ un vecteur non nul.

On note $\mathcal{P}$ le plan passant par le point $A$ et admettant le vecteur $\vec n$ comme vecteur normal.

Un point $M(x,y,z)$ appartient au plan $\mathcal{P}$ si et seulement si\[ a(x-x_A)+b(y-y_A)+c(z-z_A)=0.\]

Cette équation est appelée équation cartésienne du plan $\mathcal{P}$.


Réciproquement, si $a$, $b$, $c$ et $d$ sont quatre réels fixés, avec $(a,b,c)\neq (0,0,0)$, l'ensemble des points $M(x;y;z)$ vérifiant $ax+by+cz+d=0$ est un plan auquel le vecteur \renewcommand{\arraystretch}{1}\coorde{n}{a}{b}{c} est  normal.\end{proposition}

\begin{demonstration} Soit $M(x;y;z)$ un point de l'espace. Le point $M$ appartient au plan $\mathcal{P}$ passant par $A (x_A,y_A,z_A)$ et de vecteur normal \renewcommand{\arraystretch}{1}\coorde{n}{a}{b}{c} si et seulement si $\overrightarrow{AM}\cdot \vec n =0$. Puisque le repère que l'on considère est orthonormé,
\[\overrightarrow{AM}\cdot \vec n =  \renewcommand{\arraystretch}{1}\coordbe{x-x_A}{y-y_A}{z-z_A} \cdot \renewcommand{\arraystretch}{1}\coordbe{a}{b}{c}= a(x-x_A)+b(y-y_A)+c(z-z_A).\]
Ainsi, le point $M$ appartient au plan $\mathcal{P}$ si et seulement si $ a(x-x_A)+b(y-y_A)+c(z-z_A)=0$.

Réciproquement, $a$, $b$, $c$ et $d$ quatre réels fixés, avec $(a,b,c)\neq (0,0,0)$. On supposera par exemple que $a\neq 0$. On note alors $\vec n$ le vecteur de coordonnées $\renewcommand{\arraystretch}{1}\coordbe{a}{b}{c}$.

D'une part, l'ensemble $\mathcal{E}$ des points de l'espace vérifiant $ax+by+cz+d=0$ n'est pas vide. En effet, le point $A\left( -\frac{d}{a} ; 0 ; 0\right)$ appartient à cet ensemble. 

Soit $M(x;y;z)$ un point de l'espace et $A \in \mathcal{E}$. On a alors $ax_A+by_A+cz_A=-d$. Ainsi, 
\[ M\in \mathcal{E} \Leftrightarrow ax+by+cz+d=0 \Leftrightarrow ax+by+cz-ax_a-by_A-cz_A=0 \Leftrightarrow a(x-x_A)+b(y-y_A)+c(z-z_A)=0 \]

C'est-à-dire $M \in \mathcal{E}$ si et seulement si $\overrightarrow{AM} \cdot \vec n = 0$. L'ensemble $\mathcal{E}$ est donc un plan admettant $\vec n$ comme vecteur normal.\end{demonstration}


\begin{example}On considère le plan $\mathcal{P}$ d'équation cartésienne $3x+2y-3z+1=0$.\\ Ce plan admet le vecteur $\qquad\qquad$ comme vecteur normal.

\vskip30pt
Considérons le point $A(1;1;2)$. 

\vskip20pt

En revanche, le point $B(1;5;0)$ n'appartient pas à ce plan :\end{example}


\begin{example} Le plan $\mathcal{P}$ passant par $A(1;5;7)$ et admettant le vecteur \renewcommand{\arraystretch}{1}\coorde{n}{4}{-2}{3} admet pour équation cartésienne 

\vskip20pt
\end{example}
Il est aussi possible de raisonner comme suit : tout plan admettant le vecteur \renewcommand{\arraystretch}{1}\coorde{n}{4}{-2}{3} comme vecteur normal admet une équation cartésienne de la forme $4x-2y+3z+d=0$ pour un certain réel $d$. Pour que ce plan passe par le point $A$, il faut que les coordonnées de $A$ vérifient cette équation.

Autrement dit, $4 \times 1 -2 \times 5 + 3 \times 7 + d =0$, soit $15+d=0$ et donc $d=-15$.





\subsection{Application : intersection d'une droite et d'un plan}



\begin{example}On considère le plan $\mathcal{P}$ d'équation $2x+5y-3z+1=0$ et la droite $(d)$ de représentation paramétrique  $\left\{ \begin{array}{l}x=2-t \\ y=5+t \\ z = -2-2t \\\end{array}\right., t \in \mathbb{R}$. Soit $M(x;y;z)$ un point de l'espace.

\vskip200pt \end{example}
\newpage


\section{Projeté orthogonal}

\subsection{Projeté orthogonal sur une droite}

\begin{definition}Soit $A$ un point de l'espace et $(d)$ une droite de l'espace, dirigée par un vecteur $\vec u$.
On appelle projeté orthogonal de $A$ sur $(d)$ le point $H$ de la droite $(d)$ tel que $\overrightarrow{AH} \cdot \vec{u}=0$. En particulier,
\begin{itemize}
\item Si $A$ appartient à la droite $(d)$, ce point est son propre projeté,
\item sinon, la droite $(AH)$ est perpendiculaire à la droite $(d)$.
\end{itemize}\end{definition}

\begin{center}
\includegraphics[scale=0.3]{projort}
\end{center}

$H$ est en fait l'intersection de la droite $(d)$ et du plan $\mathcal{P}$ qui passe par le point $A$ et auquel la droite $(d)$ est normale. Si l'on connaît un point et un vecteur directeur de la droite $(d)$, il est alors possible de déterminer sa représentation paramétrique, déterminer l'équation cartésienne du plan $\mathcal{P}$, puis d'en déterminer le point d'intersection avec la méthode vue précédemment : nous déterminons ainsi le projeté orthogonal de $A$ sur $(d)$.

\begin{example}Soit $(d)$ la droite de représentation paramétrique  $\left\{ \begin{array}{l}x=1-t \\ y=3+t \\ z = 1-2t \\\end{array}\right., t \in \mathbb{R}$ et $A$ le point de coordonnées $(3;3;3)$. On vérifie facilement que le point $A$ n'appartient pas à la droite $(d)$.

Un vecteur directeur de la droite $(d)$ est le vecteur $\vec n\renewcommand{\arraystretch}{1}\begin{pmatrix}-1\\1\\-2\end{pmatrix}$. Notons alors $\mathcal{P}$ le plan passant par $A$ est donc $\vec u$ est un vecteur normal : la droite $(d)$ sera ainsi orthogonale à ce plan $\mathcal{P}$. 

Une équation cartésienne du plan $\mathcal{P}$ est donc 

\vskip30pt

Le projeté orthogonal de $A$ sur $(d)$ est le point d'intersection de $(d)$ et de $\mathcal{P}$.

\vskip100pt \end{example}
\newpage

\subsection{Projeté orthogonal sur un plan}
\begin{definition}Soit $A$ un point de l'espace, $\mathcal{P}$ un plan de l'espace et $\vec n$ un vecteur normal de $\mathcal{P}$. Le projeté orthogonal de $A$ sur $\mathcal{P}$ est le point d'intersection $H$ du plan $\mathcal{P}$ et de la droite passant par $A$ et dirigée par le vecteur $\vec n$. En particulier
\begin{itemize}
\item Si $A$ appartient au plan $\mathcal{P}$, ce point est son propre projeté
\item sinon, le vecteur $\overrightarrow{AH}$ est normal au plan $\mathcal{P}$
\end{itemize}\end{definition}

\begin{center}
\includegraphics[scale=0.3]{projort2}
\end{center}


\begin{example} On se place dans un repère orthonormé $\Oijk$. On considère le point $A$ de coordonnées $(1;3;6)$, le point $B$ de coordonnées $(1;1;1)$ et le plan $\mathcal{P}$ passant par $B$ et dirigé par \renewcommand{\arraystretch}{1}\coorde{v_1}{1}{0}{0} et \renewcommand{\arraystretch}{1}\coorde{v_2}{1}{10}{-4}. Ces deux vecteurs n'étant pas colinéaires, on définit bien ainsi un plan.

\begin{itemize}
\item Les coordonnées du vecteur $\overrightarrow{AB}$ sont $\renewcommand{\arraystretch}{1}\coorde{AB}{0}{-2}{-5}$.
\item Puisque l'on est dans un repère orthonormé, il est possible de calculer le produit scalaire à l'aide des coordonnées,
\begin{itemize}
\item $\overrightarrow{AB}\cdot \vec v_1 = 0 \times 1 + (-2) \times 0 + (-5) \times 0 = 0$ ;
\item $\overrightarrow{AB}\cdot \vec v_1 = 0 \times 2 + (-2) \times 10 + (-5) \times (-4) = 0$ ;
\end{itemize}
\item Ainsi, $B$ appartient au plan $\mathcal{P}$ et le vecteur $\overrightarrow{AB}$ est normal au plan $\mathcal{P}$. $B$ est donc le projeté orthogonal de $A$ sur $\mathcal{P}$.
\end{itemize}\end{example}

Là encore, une méthode similaire à la précédente peut être utilisée pour déterminer par le calcul les coordonnées du projeté orthogonal d'un point sur un plan. Si l'on connaît l'équation cartésienne d'un plan $\mathcal{P}$ ainsi que les coordonnées d'un point $A$ n'appartenant pas à ce plan, il est possible de déterminer les coordonnées du projeté orthogonal de $A$ sur $\mathcal{P}$. 

Pour cela, on identifie un vecteur $\vec n$ normal à $\mathcal{P}$ et on détermine une représentation paramétrique de la droite passant par $A$ et admettant $\vec n$ comme vecteur directeur : cette droite est orthogonale au plan $\mathcal{P}$. Il nous reste alors à déterminer les coordonnées du point d'intersection de cette droite et de ce plan pour obtenir le projeté orthogonal de $A$ sur $\mathcal{P}$.

\newpage
\subsection{Distance à une droite ou un plan}

\begin{definition}Soit $A$ un point de l'espace et $(d)$ une droite (ou $\mathcal{P}$ un plan). 

On appelle distance de $A$ à $(d)$ (ou à $\mathcal{P}$) la plus petite distance $AM$ pour $M$ un point de la droite $(d)$ (ou du plan $\mathcal{P}$).\end{definition}

\begin{example}Soit $A(2,1,3)$ et $(d)$ la droite de représentation paramétrique $\left\{ \begin{array}{l}x=1-t \\ y=5+t \\ z = -2-2t \\\end{array}\right., t \in \mathbb{R}$. \\Soit $M$ un point de paramètre $t$ de cette droite. On a alors
\[AM^2=\]
Développons alors cette expression. On a
\[AM^2=\]
Ainsi, $AM=\sqrt{6t^2+30t+42}$. Or, pour tout réel $t$, $6t^2+30t+42>0$. On rappelle que si $u$ est une fonction définie, dérivable et strictement positive sur $\mathbb{R}$, alors la fonction $\sqrt{u}$ est dérivable sur $\mathbb{R}$ et $(\sqrt{u})'=\dfrac{u'}{2\sqrt{u}}$. Ainsi, la fonction $f:x\mapsto \sqrt{6t^2+30t+42}$ est donc définie et dérivable sur $\mathbb{R}$ et, pour tout réel $t$, 
\[f'(t)= \]


\vskip300pt

\end{example}

\newpage

\begin{proposition}Soit $\mathcal{S}$ un plan ou une droite de l'espace.

Soit $A$ un point de l'espace et $H$ le projeté orthogonal de $A$ sur $\mathcal{S}$.

Pour tout point $K$ de l'ensemble $\mathcal{S}$ ,on a $AK \geqslant AH$ : le projeté orthogonal de $A$ sur $\mathcal{S}$ est le point de l'ensemble $\mathcal{S}$ qui est le plus proche du point $A$. La distance du point $A$ à l'ensemble $\mathcal{S}$ est alors égale à la distance $AH$.\end{proposition}

\begin{demonstration} Soit $K$ un point de $\mathcal{S}$. 

\vskip150pt

\end{demonstration}

\begin{example}On considère la droite $(d)$ et le point $A$ de l'exemple précédent. On a vu que la distance de $A$ à $(d)$ valait $\dfrac{3\sqrt{2}}{2}$ et que cette distance était atteinte pour le paramètre $t=-\dfrac{5}{2}$ dans la représentation paramétrique de $(d)$ donnée.

Ainsi, le point le plus proche du point $A$ est le point de coordonnées $(x;y;z)$ avec
\[\left\{\renewcommand{\arraystretch}{2} \begin{array}{l}x=1-\left(-\dfrac{5}{2}\right)=\dfrac{7}{2} \\ y=5+\left(-\dfrac{5}{2}\right)=\dfrac{5}{2} \\ z = -2-2 \times \left(-\dfrac{5}{2}\right)=3 \\\end{array}\right. .\]

Le point $H\left(\dfrac{7}{2};\dfrac{5}{2};3\right)$ est donc le projeté orthogonal du point $A$ sur la droite $(d)$.\end{example}





\end{document}